\paragraph{A uniform far-zero kernel remainder.}
Let $0=t_0\le\cdots\le t_N=D$ and let $q$ be continuous on
$[0,D]$, with $q(0)=1$, $q(D)=0$ and $q(x)=0$ for $x>D$.
On the $j$th closed interval assume that $q'$ has a continuous
one-sided extension $d_j$ and that $d'_j=e_j$ on its interior,
where $e_j$ is integrable. Write $E(x)=e^{-wx}$ with $w=Lz$,
$L>0$ and $\Re z>0$. Two integrations by parts give the identity
on a cell $[l,r]$
\[
 w^2\int_l^r qE-w\{q(l)E(l)-q(r)E(r)\}
   =d_j(l)E(l)-d_j(r)E(r)+\int_l^r e_j E.
\]
The terms involving $q$ telescope. Since $|E(x)|\le1$ for
$x\ge0$, the triangle inequality gives
\[
 \left|L\int_0^\infty q(x)e^{-Lzx}\,dx-\frac1z\right|
 \le \frac1{L|z|^2}
  \sum_{j=0}^{N-1}\left(
       |d_j(t_j)|+|d_j(t_{j+1})|
       +\int_{t_j}^{t_{j+1}}|e_j(x)|\,dx\right).
\]
This bound permits jumps in the first derivative. On a single
interval with $d_0(D)=0$ it has the familiar numerator
$|q'(0+)|+\int_0^D |q''(x)|\,dx$.
The half-line integral, the two integrations by parts, the
partition sum and the scaling are formalized in
\texttt{KernelRemainder.halfline\_scaled\_bound}.
The corollary \texttt{real\_far\_sum} applies this estimate to an
entire countable family of far zeros. Its conclusions are absolute
summability of the real kernel differences and the bound
$C M(1+D_0/\eta)^2$ from the stated reference denominator budget,
where $D_0$ is the displacement to the reference point.
